\documentclass[12pt]{amsart}

\usepackage{amsmath,amssymb,amsthm}
\usepackage{booktabs}
\usepackage{hyperref}
\usepackage{enumitem}

\newtheorem{theorem}{Theorem}[section]
\newtheorem{proposition}[theorem]{Proposition}
\newtheorem{lemma}[theorem]{Lemma}
\newtheorem{corollary}[theorem]{Corollary}
\theoremstyle{definition}
\newtheorem{definition}[theorem]{Definition}
\newtheorem{conjecture}[theorem]{Conjecture}
\theoremstyle{remark}
\newtheorem{remark}[theorem]{Remark}
\newtheorem{question}[theorem]{Question}
\newcommand{\leg}[2]{\left(\frac{#1}{#2}\right)}

\title[Unconditional Density Bound on ES Exceptions]{An Unconditional Density Bound on
Erd\H{o}s--Straus Exceptional Primes\\
via the Sharpness of Gateway $A = 7$}

\author{Macdara \'O Murch\'u}

\begin{document}

\begin{abstract}
Let $E(X)$ denote the number of primes $p \le X$ for which the
Erd\H{o}s--Straus equation $4/p = 1/x + 1/y + 1/z$ has no solution
in positive integers and which are not covered by any of the known
algebraic or gateway decompositions in the framework of~\cite{ES-paper}.
We prove the first unconditional quantitative bound:
\[
  E(X) \;=\; O\!\left(\frac{X}{(\log X)^{3/2}}\right),
\]
improving the trivial bound $O(X/\log X)$ by a factor of $(\log X)^{1/2}$.
The proof uses the \emph{sharpness} of gateway $A=7$
(every Case~B QR7 prime with an NQR prime factor in $N_7 = (p+7)/4$
is solved by $A=7$, proved unconditionally in~\cite{closure-paper})
to recast the exceptional set as a shifted-prime sequence sifted by the
quadratic non-residues modulo $7$, a sieve problem of dimension
$\kappa = \tfrac12$. A half-dimensional Selberg sieve with a
Bombieri--Vinogradov level of distribution then gives the bound; only an
upper bound is sought, so the parity obstruction does not intervene.
We also characterise exactly which $A$ values can be ``sharp'' in
this sense, confirm that $A=7$ is the unique sharp gateway, and
describe the character-sum barrier that prevents extending the
method to give $E(X) = O(1)$.
\end{abstract}

\maketitle

% ======================================================================
\section{Introduction}
% ======================================================================

The Erd\H{o}s--Straus conjecture (ES) asserts that for every integer
$n \ge 2$ there exist positive integers $x, y, z$ satisfying
\[
  \frac{4}{n} \;=\; \frac{1}{x} + \frac{1}{y} + \frac{1}{z}.
\]
It suffices to verify the conjecture for prime values of~$n$.
Write $E(X)$ for the number of prime counter-examples up to~$X$.
Elementary congruence arguments~\cite{ES-paper} reduce the
problem to the \emph{Case~B QR7} residue class: primes
$p \equiv 1 \pmod{24}$ with $p \equiv 1, 2, \text{ or } 4 \pmod{7}$
for which the fraction $(p+3)/4$ is $Q_3$-smooth
(all prime factors $\equiv 1 \pmod{3}$).
These primes constitute approximately $2.4\%$ of all primes,
but they resist the standard algebraic decompositions.

The \emph{gateway decomposition framework}~\cite{ES-paper} handles
each Case~B QR7 prime using candidate parameters $A$, primes
$\equiv 3 \pmod{4}$, tried in increasing order until one succeeds.
An initial implementation fixed a list of $28$ values
$A \in \{3,7,11,\ldots,239\}$ and confirmed that one of them always
succeeds for all primes $p \le 10^9$ (with $A = 239$ the worst case);
Section~\ref{sec:sharp} analyses exactly this $28$-value list.
Section~\ref{sec:computation} reports the corrected, baseline-validated
verification of all $4{,}118{,}054{,}813$ primes to $10^{11}$
(superseding an earlier run with an inverted Case~A/Case~B classifier,
see Remark~\ref{rem:integrity}), searching with no fixed candidate cap:
only $32$ distinct values of $A$ are ever required, the largest being
$A = 359$ at $p = 3{,}807{,}728{,}761$. Exactly $47$ primes below
$10^{11}$ require $A \ge 199$, and only five require $A \ge 251$; there
are no open cases.

Theoretically, the only prior bound was the trivial
\[
  E(X) \;\le\; \pi_{\mathrm{CaseB}}(X) \;=\; O\!\left(\frac{X}{\log X}\right),
\]
where $\pi_{\mathrm{CaseB}}(X)$ counts all Case~B QR7 primes up to~$X$.
Our main result is:

\begin{theorem}[Main]\label{thm:main}
$E(X) = O\bigl(X / (\log X)^{3/2}\bigr)$.
The implicit constant is effectively computable.
\end{theorem}

The key ingredient is the exact characterisation of gateway $A=7$
failure, proved in~\cite{closure-paper}:

\begin{theorem}[A=7 Sharpness,~{\cite[Theorem~5.2]{closure-paper}}]
\label{thm:a7sharp}
For every Case~B QR7 prime $p$, gateway $A=7$ succeeds if and only if
$N_7 = (p+7)/4$ has at least one prime factor that is a quadratic
non-residue modulo~$7$.
\end{theorem}

This iff characterisation is the crucial bridge: it allows us to
translate gateway failure into a multiplicative smoothness condition,
amenable to analytic methods.

\begin{remark}
Theorem~\ref{thm:a7sharp} holds for $A=7$ and \emph{no other} candidate
$A$ in the framework.  Section~\ref{sec:sharp} explains exactly why.
\end{remark}

\subsection*{Organisation}
Section~\ref{sec:setup} recalls notation and the gateway machinery.
Section~\ref{sec:bound} proves Theorem~\ref{thm:main}.
Section~\ref{sec:sharp} characterises sharp gateways.
Section~\ref{sec:charsums} describes the character-sum approach and
the barrier to unconditional finitude.
Section~\ref{sec:computation} reports the verification results to $10^{11}$.
Section~\ref{sec:open} lists open problems.

% ======================================================================
\section{Setup and Notation}\label{sec:setup}
% ======================================================================

We follow the notation of~\cite{ES-paper,closure-paper} throughout.

For a prime $A \equiv 3 \pmod{4}$, write $Q_A$ for the set of
quadratic residues modulo~$A$ (including $0$):
\[
  Q_A \;=\; \{r \in \mathbb{Z}/A\mathbb{Z} : r \equiv x^2 \pmod{A}
  \text{ for some } x\}.
\]
The non-trivial quadratic residues form a subgroup of $(\mathbb{Z}/A\mathbb{Z})^*$
of index~$2$ and order $(A-1)/2$.

\begin{definition}[$Q_A$-smooth]
A positive integer $n$ is \emph{$Q_A$-smooth} if every prime factor $q$
of~$n$ satisfies $\leg{q}{A} = +1$.
\end{definition}

For a Case~B QR7 prime~$p$ and a candidate parameter $A \equiv 3\pmod{4}$
with $4 \mid (p+A)$, set $N_A = (p+A)/4$.  The \emph{gateway target} is
\[
  t_A \;=\; {-4N_A^2} \bmod A,
\]
and the \emph{gateway set} is
\[
  S_A(p) \;=\; \bigl\{d \bmod A : d \mid N_A^2,\; \gcd(d, Ap) = 1\bigr\}.
\]
Gateway $A$ \emph{succeeds} for prime~$p$ if $t_A \in S_A(p)$.

The key lemma~\cite[Lemma~3.1]{ES-paper} asserts:
\begin{lemma}[NQR Target]\label{lem:nqr-target}
For every prime $A \equiv 3 \pmod{4}$, the target $t_A$ is a quadratic
non-residue modulo~$A$.
\end{lemma}
This follows immediately from $(-1/A) = -1$ when $A \equiv 3 \pmod{4}$.

\begin{theorem}[Failure Characterisation,~{\cite[Theorem~3.2]{closure-paper}}]
\label{thm:fc}
\begin{enumerate}[label=\textup{(\roman*)}]
  \item If $N_A$ is $Q_A$-smooth, then $S_A(p) \subseteq Q_A \cup \{0\}$,
    so $t_A \notin S_A(p)$ and gateway $A$ fails.
  \item Conversely, if gateway $A$ succeeds, then $N_A$ is not $Q_A$-smooth.
  \item Having an NQR prime factor is necessary but not sufficient
    for gateway success: $S_A(p)$ may contain NQR residues without
    containing the specific target~$t_A$.
\end{enumerate}
\end{theorem}

From (i): $\{p : N_A \text{ is } Q_A\text{-smooth}\}
          \subseteq \{p : \text{gateway } A \text{ fails}\}$.

From (ii): $\{p : \text{gateway } A \text{ succeeds}\}
           \subseteq \{p : N_A \text{ is not } Q_A\text{-smooth}\}$.

% ======================================================================
\section{The Unconditional Density Bound}\label{sec:bound}
% ======================================================================

\subsection{Reduction to a smoothness count}

\begin{lemma}\label{lem:contain}
$E(X) \le \#\{p \le X : p \text{ Case~B QR7},\; N_7 \text{ is } Q_7\text{-smooth}\}$.
\end{lemma}

\begin{proof}
If $p$ is an exceptional prime ($E(X)$ counter-example), every gateway fails.
In particular, gateway $A=7$ fails.
By Theorem~\ref{thm:a7sharp}, the exact characterisation for $A=7$,
gateway $A=7$ fails if and only if $N_7$ is $Q_7$-smooth.
Hence every exceptional prime has $N_7$ $Q_7$-smooth.
\end{proof}

\begin{remark}
The inclusion is tight: if a Case~B QR7 prime $p$ has $N_7$ $Q_7$-smooth,
then gateway $A=7$ certainly fails, but some other gateway $A \ne 7$
might still succeed for~$p$.  Thus the right-hand side of
Lemma~\ref{lem:contain} is in general strictly larger than $E(X)$.
\end{remark}

\subsection{Heuristic motivation: the density of $Q_7$-smooth integers}

\begin{lemma}[Landau--Selberg for Chebotarev sets]\label{lem:qa-density}
Let $\mathcal{P}$ be a set of primes with Dirichlet (Chebotarev) density
$\delta \in (0,1)$.  Call $n$ \emph{$\mathcal{P}$-smooth} if every prime
factor of~$n$ lies in~$\mathcal{P}$.  Then
\[
  \#\{n \le X : n \text{ is } \mathcal{P}\text{-smooth}\}
  \;\sim\; C_{\mathcal{P}}\,\frac{X}{(\log X)^{1-\delta}}
\]
for an explicit positive constant $C_{\mathcal{P}}$.
\end{lemma}

This is a classical consequence of the Selberg--Delange method applied
to the Dirichlet series
$\sum_{n \mathcal{P}\text{-smooth}} n^{-s}
= \prod_{q \in \mathcal{P}}(1-q^{-s})^{-1}$,
which has a singularity of order~$\delta$ at $s=1$;
see~\cite[Chapter~II.5]{Tenenbaum} for the full statement.

For $Q_7$-smooth integers: the set of QR primes mod~$7$ has
Chebotarev density $\delta = 1/2$ (half of all primes, by quadratic
reciprocity and Dirichlet's theorem).  Lemma~\ref{lem:qa-density} gives:
\begin{equation}\label{eq:q7-density}
  \#\{n \le X : n \text{ is } Q_7\text{-smooth}\}
  \;\sim\; C_7\,\frac{X}{(\log X)^{1/2}}.
\end{equation}

If the values $N_7 = (p+7)/4$ behaved, with respect to $Q_7$-smoothness,
like random integers of size $X/4$, then by~\eqref{eq:q7-density} a
proportion $\asymp (\log X)^{-1/2}$ of the $\asymp X/\log X$ admissible
primes would have $N_7$ smooth, predicting $E(X) \asymp X/(\log X)^{3/2}$.
We make this rigorous as an \emph{upper} bound. The integer $N_7$ is a
shifted prime, not a random integer, so the heuristic must be replaced
by a sieve. Because only an upper bound is required, the parity
obstruction does not arise, and a half-dimensional sieve with a
Bombieri--Vinogradov level of distribution suffices.

\subsection{The exceptional set as a sifted sequence}

Fix the three residue classes modulo $168 = \operatorname{lcm}(24,7)$
determined by $p \equiv 1 \pmod{24}$ and $p \equiv 1,2,4 \pmod 7$;
write $\mathcal{R} \subseteq (\mathbb{Z}/168\mathbb{Z})^*$ for this set,
so $|\mathcal{R}| = 3$. Define the sifting sequence
\[
  \mathcal{A} \;=\; \Bigl\{\, (p+7)/4 \;:\; p \le X \text{ prime},\;
  p \bmod 168 \in \mathcal{R} \,\Bigr\},
\]
and the sifting set of primes (the non-residues modulo $7$)
\[
  \mathcal{P} \;=\; \bigl\{\, \ell \text{ prime} :
  \ell \equiv 3,5,6 \pmod 7 \,\bigr\}.
\]
Write $P(z) = \prod_{\ell \in \mathcal{P},\, \ell < z} \ell$ and
\[
  S(\mathcal{A}, \mathcal{P}, z)
  \;=\; \#\{\, a \in \mathcal{A} : \gcd(a, P(z)) = 1 \,\}.
\]

\begin{lemma}[Smoothness forces survival]\label{lem:survive}
For every $z \ge 2$,
\[
  \#\{p \le X : p \bmod 168 \in \mathcal{R},\;
  N_7 \text{ is } Q_7\text{-smooth}\}
  \;\le\; S(\mathcal{A}, \mathcal{P}, z).
\]
\end{lemma}

\begin{proof}
If $N_7 = (p+7)/4$ is $Q_7$-smooth then every prime factor of $N_7$ is a
quadratic residue modulo $7$, so $N_7$ has no prime factor in
$\mathcal{P}$; in particular $\gcd(N_7, P(z)) = 1$, and $p$ contributes
the element $a = N_7$ to the survivors. Here we use that for odd
$\ell \ne 7$, $\ell \mid (p+7)/4 \iff \ell \mid (p+7)$ since
$\gcd(\ell, 4) = 1$; every $\ell \in \mathcal{P}$ is such, because
$2, 7 \notin \mathcal{P}$ (indeed $2 \equiv 3^2$ is a residue and
$7 \equiv 0 \pmod 7$).
\end{proof}

Combining Lemmas~\ref{lem:contain} and~\ref{lem:survive},
$E(X) \le S(\mathcal{A}, \mathcal{P}, z)$ for every $z \ge 2$.

\subsection{Local densities and the sieve dimension}

\begin{lemma}\label{lem:local}
For squarefree $d \mid P(z)$ write
$\mathcal{A}_d = \{a \in \mathcal{A} : d \mid a\}$.
Then $\gcd(d, 168) = 1$ and
\[
  |\mathcal{A}_d| \;=\; \frac{\omega(d)}{d}\, X_0 \;+\; r_d,
  \qquad
  \frac{\omega(d)}{d} = \frac{1}{\phi(d)},
  \qquad
  X_0 = \frac{3}{\phi(168)} \operatorname{Li}(X),
\]
where
$r_d = \sum_{i=1}^{3}\bigl(\pi(X; 168d, c_i)
- \operatorname{Li}(X)/\phi(168d)\bigr)$
for residues $c_i$ supplied by the Chinese Remainder Theorem.
The density $\omega$ is multiplicative with
$\omega(\ell)/\ell = 1/(\ell-1)$ for $\ell \in \mathcal{P}$, and the
sieve has dimension $\kappa = \tfrac12$:
\[
  V(z) \;:=\; \prod_{\substack{\ell < z \\ \ell \in \mathcal{P}}}
  \Bigl(1 - \frac{\omega(\ell)}{\ell}\Bigr)
  \;=\; \prod_{\substack{\ell < z \\ \ell \in \mathcal{P}}}
  \Bigl(1 - \frac{1}{\ell-1}\Bigr)
  \;\asymp\; (\log z)^{-1/2}.
\]
\end{lemma}

\begin{proof}
For $\ell \in \mathcal{P}$ we have $\ell \nmid 168$, and as in
Lemma~\ref{lem:survive} $\ell \mid a = (p+7)/4 \iff p \equiv -7 \pmod \ell$.
Since $\gcd(-7, \ell) = 1$ this is one class modulo $\ell$, coprime
to $\ell$. By the Chinese Remainder Theorem the conditions
``$p \bmod 168 \in \mathcal{R}$ and $p \equiv -7 \pmod d$'' select
exactly $3$ classes modulo $168d$, all coprime to $168d$ (using
$\gcd(d,168)=1$). Counting primes in these classes gives $|\mathcal{A}_d|$
with main term
$3\operatorname{Li}(X)/\phi(168d)
= \tfrac{1}{\phi(d)} \cdot \tfrac{3\operatorname{Li}(X)}{\phi(168)}$,
i.e. $\omega(d)/d = 1/\phi(d)$, multiplicative with
$\omega(\ell)/\ell = 1/(\ell-1)$.

For the dimension,
$\log(1 - \tfrac{1}{\ell-1}) = -\tfrac{1}{\ell} + O(\ell^{-2})$, so
\[
  -\log V(z)
  = \sum_{\substack{\ell < z \\ \ell \in \mathcal{P}}} \frac{1}{\ell} + O(1)
  = \tfrac12 \log\log z + O(1),
\]
by Mertens' theorem for arithmetic progressions: the residues
$3, 5, 6 \pmod 7$ comprise three of the $\phi(7) = 6$ reduced classes,
each contributing $\tfrac16 \log\log z + O(1)$ to
$\sum_{\ell < z} 1/\ell$. Exponentiating gives $V(z) \asymp (\log z)^{-1/2}$.
\end{proof}

\subsection{Level of distribution}

\begin{lemma}[Bombieri--Vinogradov input]\label{lem:bv}
For every $K > 0$, with $D = X^{2/5}$,
\[
  \sum_{\substack{d \le D \\ d \mid P(z)}}
  \mu^2(d)\, 3^{\nu(d)}\, |r_d|
  \;\ll_K\; \frac{X}{(\log X)^{K}}.
\]
\end{lemma}

\begin{proof}
By the triangle inequality, writing
$\Delta(q, a) = \pi(X; q, a) - \operatorname{Li}(X)/\phi(q)$,
\[
  |r_d| \;\le\; \sum_{i=1}^{3} |\Delta(168d, c_i)|
  \;\le\; 3 \max_{(a, 168d)=1} |\Delta(168d, a)|,
\]
and $168d \le 168 D \le X^{1/2}$ for large $X$. Cauchy--Schwarz gives
\[
  \sum_{d \le D} 3^{\nu(d)} |r_d|
  \le \Bigl(\sum_{d \le D} \frac{9^{\nu(d)}}{\phi(d)} \cdot X \log X\Bigr)^{1/2}
      \Bigl(\sum_{d \le D} |r_d|\Bigr)^{1/2},
\]
using the trivial bound $|r_d| \ll X \log X / \phi(d)$. The first factor
is $\ll X^{1/2} (\log X)^{5}$ since
$\sum_{d \le D} 9^{\nu(d)}/\phi(d) \ll (\log X)^{9}$; the second is
$\ll \bigl(X/(\log X)^{K'}\bigr)^{1/2}$ for any $K'$ by the
Bombieri--Vinogradov theorem~\cite[Thm.~17.1]{Iwaniec-Kowalski},
valid since $168 D \le X^{1/2}$. Choosing $K'$ large gives the claim.
\end{proof}

\subsection{Proof of Theorem~\ref{thm:main}}

Take $z = X^{1/5}$, so the Selberg upper-bound sieve with weights
supported on $d \le \xi = z$ has remainder over moduli $d \le D = \xi^2 = X^{2/5}$.
By Lemma~\ref{lem:local} the sequence $\mathcal{A}$ satisfies the
$\kappa$-dimensional hypothesis of Selberg's sieve with $\kappa = \tfrac12$
and density $\omega$ (note $\omega(\ell)/\ell = 1/(\ell-1) \le \tfrac12$
for $\ell \ge 3$), and by Lemma~\ref{lem:bv} it has level of distribution
$D = X^{2/5} \le X^{1/2-\epsilon}$, within Bombieri--Vinogradov range.
Selberg's sieve
\cite[Thm.~2.2 and the $\kappa$-dimensional bound of \S2.4]{Halberstam-Richert},
\cite[Ch.~6]{Iwaniec-Kowalski} yields
\[
  S(\mathcal{A}, \mathcal{P}, z)
  \;\ll\; X_0\, V(z)
  \;+\; \sum_{\substack{d \le D \\ d \mid P(z)}}
  \mu^2(d)\, 3^{\nu(d)}\, |r_d|.
\]
By Lemma~\ref{lem:local}, $V(z) \asymp (\log z)^{-1/2}
\asymp (\log X)^{-1/2}$ since $\log z = \tfrac15 \log X$, and
$X_0 \asymp X/\log X$; the remainder is $\ll_K X/(\log X)^{K}$ for any
$K$ by Lemma~\ref{lem:bv}. Hence
\[
  S(\mathcal{A}, \mathcal{P}, z)
  \;\ll\; \frac{X}{\log X} \cdot \frac{1}{(\log X)^{1/2}}
  \;=\; \frac{X}{(\log X)^{3/2}}.
\]
By Lemmas~\ref{lem:contain} and~\ref{lem:survive},
$E(X) \le S(\mathcal{A}, \mathcal{P}, z)$, proving the theorem. The implied
constant is effective, as the Bombieri--Vinogradov theorem and the
Selberg bound are effective for the moduli involved. \qed

\begin{remark}[Only an upper bound is claimed]
The argument bounds $E(X)$ from above and uses no asymptotic for
$S(\mathcal{A}, \mathcal{P}, z)$. A matching lower bound, or an asymptotic
for $\#\{p \le X : N_7 \text{ is } Q_7\text{-smooth}\}$, would meet the
parity obstruction of sieve theory and is not needed here. The exponent
$3/2$ is exactly $1 + \kappa$ with $\kappa = \tfrac12$ the relative
density of non-residues modulo $7$; the saving $(\log X)^{-1/2}$ is
intrinsic to the quadratic splitting at $A = 7$.
\end{remark}

% ======================================================================
\section{Sharp Gateways and the Uniqueness of $A = 7$}\label{sec:sharp}
% ======================================================================

\begin{definition}[Sharp gateway]
A candidate parameter $A \equiv 3 \pmod{4}$ is \emph{sharp} if,
for every Case~B QR7 prime~$p$ for which $4 \mid (p+A)$,
gateway~$A$ succeeds \emph{if and only if} $N_A$ is not $Q_A$-smooth;
and moreover the characterisation is \emph{non-vacuous}: there exist
Case~B QR7 primes $p$ for which $N_A$ is not $Q_A$-smooth
(so gateway~$A$ sometimes succeeds on Case~B QR7 primes).
\end{definition}

\begin{remark}
The non-vacuity clause excludes $A = 3$.
By definition, Case~B QR7 primes satisfy $(p+3)/4$ is $Q_3$-smooth,
so $N_3 = (p+3)/4$ is always $Q_3$-smooth and gateway~$A=3$
always fails on Case~B QR7 primes.
The iff holds vacuously but gives no density improvement.
\end{remark}

By Theorem~\ref{thm:fc}(i)(ii), every gateway satisfies the forward
implication ($Q_A$-smooth $\Rightarrow$ fails) and the converse direction
at the level of necessary conditions (success $\Rightarrow$ not smooth).
A sharp gateway additionally satisfies the reverse: not smooth $\Rightarrow$
succeeds.

The proof of Theorem~\ref{thm:a7sharp} in~\cite{closure-paper} proceeds by:
\begin{enumerate}
  \item Identifying the \emph{forced} QR factor $g = 2$ of $N_7$:
    since $p \equiv 1 \pmod{8}$, we have $8 \mid (p+7)$, hence $2 \mid N_7$.
  \item Observing that $2$ \emph{generates} $Q_7 = \{1,2,4\}$:
    the powers $2^0, 2^1, 2^2$ cycle through all of $Q_7 \pmod{7}$.
  \item Showing that if $q \mid N_7$ with $\leg{q}{7} = -1$, then
    $q \cdot 2^j \bmod 7$ hits every NQR modulo~$7$ as $j \in \{0,1,2\}$,
    including the target $t_7$.
  \item Verifying the integrality condition: $d = q \cdot 2^{2k}$
    divides $N_7^2$ (since $v_2(N_7^2) = 2v_2(N_7) \ge 2$),
    and $\gcd(d, 7p) = 1$ by Lemma~3.2 of~\cite{ES-paper}.
\end{enumerate}

\begin{theorem}[Uniqueness of $A=7$]\label{thm:unique}
Among the $28$ candidate $A$ values in the framework,
$A = 7$ is the only sharp gateway.
\end{theorem}

\begin{proof}
Sharpness requires a prime $g$ that is \emph{forced} to divide $N_A$ for
every Case~B QR7 prime with $4 \mid (p+A)$, such that
(a) $g$ is a QR modulo~$A$, (b) $g$ generates $Q_A$ under multiplication
modulo~$A$, and (c) the guaranteed valuation suffices:
$2 v_g(N_A) \ge |Q_A| - 1$, so that $\{g^0, g^1, \ldots, g^{2v_g}\}$
realises all of $Q_A$ as divisor residues of $N_A^2$.

\textit{Which primes can be forced.}
Case~B QR7 fixes $p \equiv 1 \pmod 8$ and $p \equiv 1 \pmod 3$, but leaves
$p \bmod \ell$ free for every prime $\ell \ge 5$. Hence the only primes
that can divide $N_A = (p+A)/4$ for \emph{all} such $p$ are $2$ and $3$,
and a direct check of $4 N_A = p + A$ gives the rule
\[
  2 \mid N_A \iff A \equiv 7 \pmod 8,
  \qquad
  3 \mid N_A \iff A \equiv 2 \pmod 3.
\]
A candidate with neither $A \equiv 7 \pmod 8$ nor $A \equiv 2 \pmod 3$ has
no forced prime factor and cannot be sharp.

We check each candidate:

\textbf{$A = 7$:}
$7 \equiv 7 \pmod 8$, so $g = 2$ is forced; $7 \equiv 1 \pmod 3$, so $3$
is not. $(2/7) = +1$, and $\operatorname{ord}_7(2) = 3 = |Q_7|$, so $2$
generates $Q_7$. The guaranteed $v_2(N_7) \ge 1$ gives
$2 v_2(N_7) \ge 2 = |Q_7| - 1$. \textbf{Sharp.}

\textbf{$A = 11$:}
$11 \equiv 2 \pmod 3$, so $g = 3$ is forced ($3 \mid N_{11}$);
$11 \equiv 3 \pmod 8$, so $2$ is not.
$(3/11) = +1$ (QR, since $3 \in \{1,3,4,5,9\}$).
Order of~$3$ in $(\mathbb{Z}/11\mathbb{Z})^*$ is~$5 = |Q_{11}|$, so $3$
generates $Q_{11}$.
However, the minimum forced multiplicity is $v_3(N_{11}) \ge 1$, giving
$2v_3(N_{11}) \ge 2$, but we need $\ge |Q_{11}|-1 = 4$.
For $p \equiv 1, 4 \pmod{9}$ (two-thirds of cases) one has $v_3(N_{11}) = 1$,
so the required divisors $3^3$ and $3^4$ do not appear in $N_{11}^2$
and the target $t_{11}$ need not be hit.

\emph{Concrete counterexample.}
Take $p = 193$ (Case~B QR7: $193 \equiv 1 \pmod{24}$,
$193 \equiv 4 \pmod{7}$, $(p+3)/4 = 49 = 7^2$ is $Q_3$-smooth).
Then $N_{11} = (193+11)/4 = 51 = 3 \times 17$.
The factor $17$ satisfies $(17/11) = -1$ (NQR: $6 \notin Q_{11}$),
so $N_{11}$ has an NQR prime factor.
Yet $t_{11} = (-4 \cdot 51^2) \bmod 11 = 2$, while
$S_{11}(193) = \{1,3,5,6,7,9,10\}$ (all residues of divisors of
$51^2 = 3^2 \cdot 17^2$ modulo~$11$), and $2 \notin S_{11}(193)$.
Gateway $A = 11$ fails for $p = 193$ despite $N_{11}$ having an NQR factor.
\textbf{Not sharp.}

\textbf{$A = 19$:}
$19 \equiv 3 \pmod 8$ and $19 \equiv 1 \pmod 3$, so \emph{neither} $2$ nor
$3$ is forced into $N_{19}$. No prime is forced, so condition for
sharpness fails at the first step. (For instance $p = 193$ is Case~B QR7
with $N_{19} = 53$, prime and coprime to $2$ and $3$.)
\textbf{Not sharp.}

\textbf{$A \ge 23$ (remaining $25$ candidates):}
A systematic check of all $28$ candidate values confirms that for every
$A \in \{23, 31, 43, \ldots, 239\}$, either
(a) no forced prime ($A \not\equiv 7 \pmod 8$ and $A \not\equiv 2 \pmod 3$,
e.g. $A = 43, 67, 79, \ldots$), or
(b) a forced prime exists but the guaranteed valuation satisfies
$2 v_g(N_A) < |Q_A| - 1$.
For $A = 23$ ($\equiv 7 \pmod 8$ and $\equiv 2 \pmod 3$) both $2$ and $3$
are forced, but $|Q_{23}| = 11$ requires $2 v_g \ge 10$ while the
guaranteed valuations are $1$ each. The growing gap $|Q_A| - 1 \ge 10$
against bounded forced valuations rules out every $A \ge 23$.
\textbf{None sharp.}

The computation was verified by testing all $28$ candidates directly,
confirming $A = 7$ is the unique potentially sharp gateway.
\end{proof}

\begin{remark}
The special arithmetic of $A = 7$:
$|Q_7| = 3$ is the smallest possible, and $|Q_7| - 1 = 2$ equals
$2 \cdot \min(v_2(N_7)) = 2 \cdot 1$.  No larger prime $A$ can match
this coincidence.
\end{remark}

% ======================================================================
\section{Character Sums and the Barrier to Finitude}\label{sec:charsums}
% ======================================================================

The density bound of Theorem~\ref{thm:main} does not imply finiteness of $E$:
$X/(\log X)^{3/2} \to \infty$ as $X \to \infty$.
We describe the character-sum approach that would be needed to
obtain $E(X) = O(1)$.

\subsection{Character sum formula for gateway solutions}

The number of gateway solutions for parameter $A$ is:
\begin{equation}\label{eq:N-formula}
  \mathcal{N}_A(p) \;=\;
  \frac{1}{\varphi(A)} \sum_{\chi \bmod A} \bar\chi(t_A)
  \sum_{\substack{d \mid N_A^2 \\ \gcd(d,Ap)=1}} \chi(d).
\end{equation}

Separating the trivial character $\chi_0$ and the Legendre symbol $\chi_2$:

\begin{align}\label{eq:split}
  \mathcal{N}_A(p) \;=\;&
    \frac{\tau'(N_A^2) - \tau'(M_A^2)}{\varphi(A)}
    \;+\; E_A(p),
\end{align}
where $\tau'(n)$ counts the divisors of~$n$ coprime to~$A$,
$M_A$ is the $Q_A$-smooth part of~$N_A$, and
$E_A(p)$ involves only characters of order~$\ge 3$.

\begin{remark}
The restriction in $\tau'$ is inherited from $\gcd(d, Ap) = 1$
in~\eqref{eq:N-formula} and is vacuous here: $4 N_A = p + A$, so
$A \mid N_A$ would force $A \mid p$, impossible for $A < p$ prime.
Hence $\tau'(N_A^2) = \tau(N_A^2)$ throughout, and likewise for $M_A$.
The prime is kept because the identity $\sum_{d \mid N_A^2,\, (d,A) = 1}
\chi_2(d) = \tau'(M_A^2)$, which produces the second term, is stated for
the coprime-restricted sum.
\end{remark}

The \emph{main term} $\tau'(N_A^2) - \tau'(M_A^2)$ is positive if and
only if $N_A$ is not $Q_A$-smooth (i.e., $M_A < N_A$).
Gateway $A$ fails only when $\mathcal{N}_A(p) = 0$,
which requires $E_A(p) \le -(\tau'(N_A^2) - \tau'(M_A^2))/\varphi(A)$.

\subsection{Why the character sum barrier is hard}

For $A = 7$: $\varphi(7) = 6$ has characters of orders $1, 2, 3, 6$.
The order-$3$ characters are the ``next-order'' corrections.
Bounding $|E_7(p)|$ requires nontrivial estimates for twisted divisor sums
$\sum_{d \mid N_7^2} \chi_3(d)$.

For general~$A$: bounding $\sum_{p \le X} |\mathcal{N}_A(p)|^{-\epsilon}$
(which would suffice for finitude under GRH-type hypotheses) is
an open problem.

\subsection{What additional hypotheses give}

\begin{theorem}[Conditional,~{\cite[Theorem~7.2]{closure-paper}}]
Under Conjecture~7.1 of~\cite{closure-paper}
(joint $Q_A$-smoothness equidistribution, weaker than Elliott--Halberstam),
the Erd\H{o}s--Straus conjecture holds for all sufficiently large primes.
\end{theorem}

The gap between Theorem~\ref{thm:main} and finiteness corresponds to
the gap between:
\begin{itemize}
  \item $O(X/(\log X)^{3/2})$ (what smoothness density gives), and
  \item $O(1)$ (what finiteness requires).
\end{itemize}
To close this gap unconditionally, one would need to show that
the character sum error $E_A(p)$ is never large enough to cancel
the main term across all 28 gateways simultaneously.

\subsection{Why more gateways do not trivially improve the exponent}

It is tempting to improve Theorem~\ref{thm:main} by using all $28$
candidates at once. The naive heuristic runs as follows. For primes
$q > 239$ each $q$ divides at most one $N_{A_i}$ (since
$N_{A_i} - N_{A_j} = (A_i - A_j)/4 \le 58 < q$ in the $28$-candidate set),
and by Chebotarev the Frobenius class of $q$ at each gateway $A_i$ is
independent. The probability that every large prime factor of $N_{A_i}$
is QR mod $A_i$ simultaneously for $i = 1, \ldots, k$ would then decay
as $(\log X)^{-k/2}$, suggesting $E(X) = O(X/(\log X)^{1+k/2})$.

This heuristic is \emph{not rigorous}, and the obstruction is structural.
The bound of Theorem~\ref{thm:main} uses the \emph{sharpness} of $A = 7$:
the exceptional set is contained in $\{p : N_7 \text{ is } Q_7\text{-smooth}\}$
because for $A = 7$, gateway failure is \emph{equivalent} to
$Q_7$-smoothness. For every other $A$ the equivalence fails
(Theorem~\ref{thm:fc}(iii)): a prime can fail gateway $A_i$ while
$N_{A_i}$ is \emph{not} $Q_{A_i}$-smooth. Consequently the exceptional
set is \emph{not} contained in the joint-smoothness set
$\bigcap_i \{N_{A_i}\text{ is }Q_{A_i}\text{-smooth}\}$, and bounding the
latter by a Selberg sieve gives \emph{no} upper bound on $E(X)$. Indeed
$\bigcap_i \{N_{A_i}\text{ smooth}\} \subseteq E$ runs the wrong way,
a lower bound on $E$ if anything.

A genuine improvement therefore requires controlling the
\emph{character-sum} failure events $\{\mathcal{N}_{A_i}(p) = 0\}$ of
\eqref{eq:N-formula} for $A_i \ne 7$, not merely smoothness. We record
the expected truth as a conjecture, with the caveat that it is not a
sieve statement:
\begin{conjecture}\label{conj:stronger}
$E(X) = O\bigl(X/(\log X)^{1+\delta}\bigr)$ for every $\delta > 0$;
conjecturally $E(X) = O(1)$.
\end{conjecture}
A proof would need joint control of the divisor character sums across the
candidate gateways, the analytic gap described in~\cite{closure-paper};
the smoothness method of Theorem~\ref{thm:main} is intrinsically limited
to the single sharp gateway $A = 7$ and the exponent $3/2$.

% ======================================================================
\section{Verification to $10^{11}$}\label{sec:computation}
% ======================================================================

\subsection{Method}

The reported numbers come from the corrected, baseline-validated
pipeline (\texttt{sessions/session20\_corrected\_10B.py} in the software
repository), which:
\begin{enumerate}
  \item generates primes with a CPU segmented sieve (segments of $10^8$);
  \item classifies each prime into the algebraic cases of~\cite{ES-paper}
    or the Case~B QR7 residual class, using the corrected Case~A/Case~B
    test (Remark~\ref{rem:integrity});
  \item for each Case~B QR7 prime, searches candidate values
    $A \equiv 3 \pmod 4$ in increasing order --- with no fixed upper cap,
    up to $A < 10^4$ --- factorising $N_A$ and testing the
    divisor-residue condition of Section~\ref{sec:setup}, recording the
    first successful~$A$;
  \item runs on a $16$-worker CPU pool with resumable, per-segment
    checkpointing, validating against three independent baselines (the
    full $A$-distribution at $10^6$, the residual count at $10^7$, and
    the residual count and $\max A$ at $10^9$) before extending the
    range.
\end{enumerate}
An earlier GPU-accelerated implementation (CuPy, fixed candidate lists
of $28$ then $87$ values) first located the $A = 359$ record described
in Section~\ref{subsec:a359}; it has since been superseded by the
pipeline above, which is the source of every number in this section.

\begin{remark}[Data integrity]\label{rem:integrity}
An earlier run of the $10^{10}$ stage (March 2026) inverted the
Case~A/Case~B test, skipping every residual-class prime as already
proven and instead gateway-searching the complementary (easy) class.
The resulting figures --- a residual count of $18{,}189{,}229$ at
$10^{10}$, $\max A = 251$, a $25$-value $A$-dictionary --- described the
wrong set of primes and are withdrawn; they do not appear in this paper.
The corrected pipeline described above reproduces the correct-classifier
baselines at $10^6$, $10^7$, and $10^9$ before extending the range, and
every reported solution is verified by direct evaluation of
$4/p = 1/x + 1/y + 1/z$.
\end{remark}

\subsection{Results to $10^{11}$}

\begin{center}
\begin{tabular}{lrrr}
\toprule
Limit & Primes & Case~B QR7 & Max $A$ required \\
\midrule
$10^6$  &    78{,}498  &    2{,}269 & 79 \\
$10^7$  &   664{,}579  &   18{,}019 & 167 \\
$10^8$  & 5{,}761{,}455 &  146{,}227 & 239 \\
$10^9$  & 50{,}847{,}534 & 1{,}212{,}383 & 239 \\
$10^{10}$ & 455{,}052{,}511 & 10{,}246{,}512 & \textbf{359} (see \S\ref{subsec:a359}) \\
$10^{11}$ & 4{,}118{,}054{,}813 & 88{,}088{,}687 & \textbf{359} (unchanged from $10^{10}$) \\
\bottomrule
\end{tabular}
\end{center}

\begin{theorem}[Computational Verification]\label{thm:gpu}
The Erd\H{o}s--Straus conjecture holds for all primes $p \le 10^{11}$,
with maximum gateway parameter $A \le 359$ and zero open cases.
Exactly $47$ primes below $10^{11}$ require $A \ge 199$ ($12$ of them
below $10^{10}$), and only five require $A \ge 251$:
$p = 3{,}807{,}728{,}761$ ($A = 359$),
$p = 31{,}048{,}497{,}481$ ($A = 347$),
$p = 6{,}372{,}847{,}849$ ($A = 311$),
$p = 26{,}147{,}464{,}729$ ($A = 251$), and
$p = 94{,}604{,}730{,}049$ ($A = 251$).
The record $A = 359$ is attained below $10^{10}$ and is not exceeded
anywhere in $(10^{10}, 10^{11}]$.
\end{theorem}

\subsection{Anatomy of the hardest prime}\label{subsec:a359}

The unique prime below $10^{11}$ requiring $A = 359$ is
\[
  p \;=\; 3{,}807{,}728{,}761.
\]

\medskip
\noindent\textit{Verification that $p$ is Case~B QR7.}
$p \equiv 1 \pmod{24}$, $p \equiv 4 \pmod 7$,
and $(p+3)/4 = 951{,}932{,}191 = 7 \times 73 \times 1{,}862{,}881$,
all three factors $\equiv 1 \pmod 3$, confirming $Q_3$-smoothness.

\medskip
\noindent\textit{Solution via $A = 359$.}
$N_{359} = (p + 359)/4 = 951{,}932{,}280
        = 2^3 \cdot 3 \cdot 5 \cdot 13 \cdot 23 \cdot 43 \cdot 617.$
The factors $13, 43, 617$ are all NQR modulo~$359$.
The divisor $d = 1935 = 3^2 \times 5 \times 43$ divides $N_{359}^2$,
$\gcd(d, 359 p) = 1$, and
$1935 \bmod 359 = 140 = t_{359}$, so gateway $A = 359$ succeeds.
This claim is independently reproduced by
\texttt{verify\_hardest\_prime.py} in the software repository.

\medskip
\noindent
This prime is the record-holder over the full $10^{11}$ range: the
largest gateway parameter required for any prime $p \le 10^{11}$.
Of the four other primes requiring $A \ge 251$ (Theorem~\ref{thm:gpu}),
one --- $p = 6{,}372{,}847{,}849$ ($A = 311$) --- also lies below
$10^{10}$, alongside the $A = 359$ record; the remaining three lie in
$(10^{10}, 10^{11}]$. None exceeds $A = 359$.

\begin{remark}
The Case~B QR7 residual class has relative density $2.14\%$ at $10^{11}$
(down from $2.89\%$ at $10^6$), consistent with the
relative-density-zero prediction for this class in~\cite{ES-paper}.
The fraction requiring large $A$ shrinks even faster: only $47$ of the
$88{,}088{,}687$ Case~B QR7 primes below $10^{11}$ need $A \ge 199$.
\end{remark}

% ======================================================================
\section{Open Problems}\label{sec:open}
% ======================================================================

\begin{question}
Improve the bound to $E(X) = O(X/(\log X)^{3/2+\epsilon})$
for some $\epsilon > 0$ using properties of gateways beyond $A=7$.
\end{question}

\begin{question}
Improve the exponent toward Conjecture~\ref{conj:stronger}
($E(X) = O(X/(\log X)^{1+\delta})$ for all $\delta > 0$) by controlling
the divisor character sums $\mathcal{N}_{A_i}(p)$ at gateways $A_i \ne 7$,
not merely smoothness. As noted in \S\ref{sec:charsums}, the joint
$Q_A$-smoothness set does not contain the exceptional set, so the
single-gateway smoothness method caps at the exponent $3/2$.
\end{question}

\begin{question}
Extend the verification beyond $10^{11}$. The corrected pipeline
(Section~\ref{sec:computation}) confirms $A = 359$ as the hardest case
for $p \le 10^{11}$, unchanged since $10^{10}$; extending further would
determine whether larger gateway parameters ever appear. The maximum
$A$ encountered for any prime $p \le X$ grows slowly with $X$
(log-like); bounding this growth rigorously is an open problem (the
\emph{bounded-$A$ conjecture} of~\cite{ES-paper}).
\end{question}

\begin{question}
Characterise the set of primes $p$ for which $N_7$ is $Q_7$-smooth.
By Theorem~\ref{thm:main}, this set has counting function
$O(X/(\log X)^{3/2})$, but no explicit family is known.
These are precisely the primes for which the A=7 gateway fails.
\end{question}

\begin{question}
Is there an $A \ne 7$ for which an exact characterisation of gateway
failure holds for a \emph{positive-density} subset of Case~B QR7 primes?
Theorem~\ref{thm:unique} rules out universal sharpness, but
a density-one subset might still be achievable for specific $A$.
\end{question}

\begin{question}
Can the character sum formula~\eqref{eq:N-formula} be combined with
the large sieve to give an unconditional proof that $\mathcal{N}_A(p) > 0$
for all Case~B QR7 primes $p \ge P_0$ (for explicit $P_0$)?
This would require bounding $E_A(p)$ in terms of the main term
and a zero-free region for $L(s, \chi)$ for characters $\chi$ mod~$A$.
\end{question}

% ======================================================================
\section*{Acknowledgements}
% ======================================================================

This work uses techniques developed in~\cite{ES-paper,closure-paper}.
An initial GPU-accelerated search (CuPy, NVIDIA Quadro RTX~5000) first
located the $A = 359$ record of Section~\ref{subsec:a359}; the results
reported here come from the corrected CPU pipeline of
Section~\ref{sec:computation}.

\begin{thebibliography}{99}

\bibitem{ES-paper}
M.~\'O~Murch\'u,
\textit{Bounded Gateway Parameters for the
Erd\H{o}s--Straus Conjecture},
preprint, 2026.

\bibitem{closure-paper}
M.~\'O~Murch\'u,
\textit{Toward Closure of the Gateway Gap:
NQR Targets, Subgroup Failure, and the Multi-$A$ Strategy},
preprint, 2026.

\bibitem{Halberstam-Richert}
H.~Halberstam and H.-E.~Richert,
\textit{Sieve Methods},
London Mathematical Society Monographs~4, Academic Press, 1974.

\bibitem{Iwaniec-Kowalski}
H.~Iwaniec and E.~Kowalski,
\textit{Analytic Number Theory},
AMS Colloquium Publications, Vol.~53, 2004.

\bibitem{Tenenbaum}
G.~Tenenbaum,
\textit{Introduction to Analytic and Probabilistic Number Theory},
Cambridge Studies in Advanced Mathematics~46, 1995.

\end{thebibliography}

\end{document}
